\section{Real scalar extension of the same original Picard group}
\label{sec:native-real-picard}

\begin{theorem}[Actual original real line classes and finite sums {\cite[\S1.2]{BJ22}}]
\label{thm:native-real-picard}
\lean{EnvelopeProperty.NativeRealPicard.Space, EnvelopeProperty.NativeRealPicard.picardClass, EnvelopeProperty.NativeRealPicard.lineClass, EnvelopeProperty.NativeRealPicard.span_picardClass, EnvelopeProperty.NativeRealPicard.span_lineClass, EnvelopeProperty.NativeRealPicard.exists_finite_line_sum}
\leanok
\uses{thm:fc-original-projective-generated-line-quotients}
For every original scheme $X$, take the real scalar extension
$\mathbf R\otimes_{\mathbf Z}\operatorname{Pic}(X)$ of its
same actual native Picard group. The class of each actual original
invertible line is represented by its same original native Picard
class. Every element of this real Picard space is a finite real
linear combination of such actual original invertible-line classes.
No abstract class-realization, finite-dimensionality, or supplied
finite-line-family certificate is assumed.
\end{theorem}

\begin{proof}
The native Picard group is the group of units in the isomorphism
quotient of the original module category. Pass to its actual
associated additive group and tensor it over the integers with
the real numbers. Send an original Picard class $a$ to $1\otimes a$,
and an original invertible line to $1\otimes[L]$. The real span
of these original Picard classes is the whole tensor product:
zero and sums are in the span, and every pure tensor $r\otimes a$
is $r(1\otimes a)$. Tensor-product induction proves the claim
for every actual tensor, without assuming finite dimension.
Every native Picard class has an actual original invertible-line
representative, obtained from its same native isomorphism-quotient
representative and the genuine unit property. Thus the real span
of actual original line classes is likewise the entire tensor
product. The actual finite-support linear-combination description
of a span supplies a finite set of these same original invertible
lines and real coefficients summing to the specified original
element. This is a genuine existence theorem for the original
lines, not a new class-realization assumption. Numerical
quotienting and topology are subsequent constructions.
\end{proof}
